\documentclass{report}
\usepackage{amsmath,amssymb}
\setlength{\parindent}{0mm}
\setlength{\parskip}{1em}
\begin{document}
\begin{center}
\rule{6in}{1pt} \
{\large Kevin Deweese \\
{\bf An Empirical Comparison of Graph Laplacian Solvers}}

Department of Computer Science \\ UC Santa Barbara \\ Santa Barbara \\ CA 93106
\\
{\tt kdeweese@cs.ucsb.edu}\\
Erik G. Boman\\
John R. Gilbert\end{center}

Solving Laplacian linear systems is an important task in a variety of
practical and theoretical applications. This problem is known to have
solutions that perform in linear times polylogarithmic work in theory,
but these algorithms are difficult to implement in practice. We examine
existing solution techniques in order to determine the best methods
currently available and for which types of problems are they useful. We
perform timing experiments using a variety of solvers on a variety of
problems and present our results.

The Laplacian matrix of a weighted, undirected graph $G$ is defined as
$L_G=D_G-A_G$, where $D_G$ is the diagonal matrix containing the sum of
incident edge weights and $A_G$ is the weighted adjacency matrix. $L_G$
is a symmetric, positive semidefinite, diagonally dominant M-matrix, with
a nullspace
containing the constant vector. Solving linear systems on the Laplacians
of structured graphs, such as two and three dimensional meshes, has long
been important to application domains such as finite element analysis and
image processing. More recently, solving linear systems on the Laplacians
of large graphs without mesh-like structure has emerged as an important
computational task in network analysis with applications to graph
sparsification and spectral clustering.

To our knowledge there is not a comprehensive experimental study of
available solvers. Our goal is to provide an initial understanding of
available solvers, and to determine where ideas from more theoretical
solvers could be leveraged to make improvements in practice. Towards this
goal we must select an implementation of existing solvers, gather a set
of relevant test problems, and select a set of performance metrics for
comparing solver performance. We choose to use the Trilinos software
package and its implementation of preconditioned conjugate gradient and
multilevel methods. We use a variety of mesh-like and irregular graphs
from the University of Florida sparse matrix collection, a set of
generated Block Two-Level Erd\"{o}s R\'{e}nyi graphs, and a set of graphs
generated from image segmentation problems. For each solver we examine
number of iterations, setup time, solve time, and memory usage.

We examine each solver's performance on each test set and present our
results. Of interest, we find that relative solver behavior is consistent
within each test set, but varies between the test sets.
Cheap to apply preconditioners such as Jacobi scaling work surprisingly
well on graphs with irregular degree distribution. We hope to update
these results with new test data, solvers, and analysis through community
feedback.


\end{document}
